% 3. 超放射の条件
\section{超放射: 回転から波を増幅する条件}\label{sec:superradiance}

\subsection{第一法則と面積定理から導く}

BH に、振動数 $\omega$、方位量子数 $m$ の波が当たるとする。波は $e^{-i\omega t+im\phi}$ の形で、量子 1 個はエネルギー $\hbar\omega$ と、BH の回転軸まわりの角運動量 $m\hbar$ をもつ。BH がこの波の量子を $\delta N$ 個吸収すると(放出なら $\delta N<0$)、
\begin{equation}
\delta M=\hbar\omega\,\delta N,\qquad \delta J=m\hbar\,\delta N.
\end{equation}
第一法則 \eqref{eq:firstlaw} に代入すると
\begin{equation}
\frac{\kappa}{8\pi}\,\delta A=\delta M-\Omega_H\delta J=\hbar\,\delta N\,(\omega-m\Omega_H).
\label{eq:sr_firstlaw}
\end{equation}
面積定理 $\delta A\ge0$ と $\kappa\ge0$ から、右辺は 0 以上でなければならない。
\begin{equation}
\delta N\,(\omega-m\Omega_H)\ge0.
\end{equation}
\begin{itemize}
\item $\omega>m\Omega_H$ のとき: $\delta N\ge0$。BH は波を吸収できる(ふつうの吸収)。
\item $\omega<m\Omega_H$ のとき: $\delta N\le0$。BH はこの波を吸収できず、量子を放出するしかない。入ってきた波は、より大きな振幅で跳ね返される。
\end{itemize}
後者を超放射と呼ぶ。条件は
\begin{equation}
\boxed{\;0<\omega<m\Omega_H\;}
\label{eq:sr_condition}
\end{equation}
である。このとき BH は質量 $\hbar\omega$ と角運動量 $m\hbar$ を 1 量子ごとに失う。失う質量は回転のエネルギー $J^2/(4M_\mathrm{irr}^2)$ から来るので、既約質量(面積)は減らない。

\paragraph{直観}
回転する物体の上に乗った観測者には、波の振動数は $\omega'=\omega-m\Omega$ に見える(角度 $\phi$ が $\phi'=\phi-\Omega t$ に変わるので、$e^{-i\omega t+im\phi}=e^{-i(\omega-m\Omega)t+im\phi'}$)。$\omega<m\Omega$ では $\omega'<0$ になり、回転系での「吸収」が、静止系では「放出」に見える。回転する円柱が電磁波を増幅する効果(Zel'dovich)と同じ仕組みである。

\subsection{質量のある粒子: 増幅が繰り返される}

質量 $\mu$ の粒子の波は、振動数が $\omega<\mu$ なら無限遠に逃げられず、BH のまわりに束縛される(水素原子の束縛状態と同じ。\S\ref{sec:atom})。束縛された波は BH に何度も当たり、超放射の条件を満たしていれば、当たるたびに増幅される。その結果、波の振幅は指数関数的に育つ。質量が「鏡」の役をする「ブラックホール爆弾」である。

束縛状態では $\omega\simeq\mu$ なので、超放射の条件 \eqref{eq:sr_condition} は $\mu<m\Omega_H$ になる。両辺に $M$ をかけ、式 \eqref{eq:OmegaH} と $\alpha=M\mu$($G=\hbar=c=1$)を使うと
\begin{equation}
\alpha<m\,\Omega_HM=\frac{m\chi}{2\rp}=\frac{m\chi}{2(1+\sqrt{1-\chi^2})}.
\label{eq:sr_alpha}
\end{equation}
右辺は $\chi=1$ で最大値 $m/2$ をとるので、$m=1$ の状態が育つには少なくとも $\alpha<1/2$ が要る。

\subsection{臨界スピン}\label{sec:chic}

BH が回転を失っていくと、やがて式 \eqref{eq:sr_alpha} の等号が成り立ち、超放射が止まる。そのときのスピン $\chi_c$ を求める。
\begin{equation}
m\chi_c=2\alpha\bigl(1+\sqrt{1-\chi_c^2}\bigr).
\end{equation}
根号を一方に集める。
\begin{equation}
m\chi_c-2\alpha=2\alpha\sqrt{1-\chi_c^2}.
\label{eq:chic_step}
\end{equation}
右辺は 0 以上なので、$m\chi_c\ge2\alpha$ が必要である。両辺を 2 乗すると
\begin{align}
m^2\chi_c^2-4m\alpha\chi_c+4\alpha^2&=4\alpha^2(1-\chi_c^2),\\
(m^2+4\alpha^2)\chi_c^2-4m\alpha\chi_c&=0.
\end{align}
$\chi_c\ne0$ の解は
\begin{equation}
\boxed{\;\chi_c=\frac{4m\alpha}{m^2+4\alpha^2}\;}
\label{eq:chic}
\end{equation}
である。これが $m\chi_c\ge2\alpha$ を満たすことは、$m\chi_c-2\alpha=2\alpha(m^2-4\alpha^2)/(m^2+4\alpha^2)\ge0$($\alpha\le m/2$ のとき)で確かめられる。

\begin{table}[H]
\centering\small
\caption{臨界スピン $\chi_c$(式 \ref{eq:chic})。$\alpha$ が大きいほど、BH は多くの回転を残す。}
\label{tab:chic}
\begin{tabular}{@{}lcccccc@{}}
\toprule
$\alpha$ & 0.02 & 0.05 & 0.08 & 0.10 & 0.15 & 0.20 \\
\midrule
$m=1$ & 0.080 & 0.198 & 0.312 & 0.385 & 0.550 & 0.690 \\
$m=2$ & 0.040 & 0.100 & 0.159 & 0.198 & 0.293 & 0.385 \\
\bottomrule
\end{tabular}
\end{table}

\begin{caution}
正確には、束縛状態の振動数は $\omega=E_{n\ell m}$ で、$\mu$ より少し小さい($E\simeq\mu(1-\alpha^2/2n^2)$、\S\ref{sec:atom})。研究のコードでは $m\Omega_H=E_{n\ell m}$ を数値的に解いている(\texttt{si.chi\_threshold\_211})。$\alpha=0.104$ での差は $\chi_c$ で $5\times10^{-4}$ ほどで、式 \eqref{eq:chic} で十分よい。
\end{caution}

\begin{summary}
\begin{itemize}
\item 面積定理と第一法則から、$\omega<m\Omega_H$ の波は BH に吸収されず増幅される(超放射)。
\item 質量のある粒子の束縛状態($\omega\simeq\mu$)は繰り返し増幅されて指数関数的に育つ。
\item 成長は BH のスピンが臨界値 $\chi_c=4m\alpha/(m^2+4\alpha^2)$ まで下がると止まる。
\end{itemize}
\end{summary}
